\section{Población, muestra y unidad de análisis}
\label{sec:unidad}

\subsection{Delimitación de la población}

La población objetivo del estudio corresponde a las estaciones de chancado
primario giratorio y de mandíbula de la gran minería a cielo abierto,
caracterizadas por alimentación discontinua mediante camiones de extracción de
gran tonelaje y por sistemas SCADA de adquisición de datos telemétricos con
frecuencias de muestreo estándar del orden de \SIrange{1}{10}{\hertz}.

\subsection{Unidad de análisis y vector de estado dinámico}

La unidad de análisis es el subsistema ciberfísico de recepción y descarga —la
tolva de alimentación del chancador— modelado computacionalmente mediante el
vector de estado dinámico bidimensional en tiempo discreto:

\begin{equation}
\mathbf{x}_{k} = \begin{bmatrix} V_{k} \\ \dot{V}_{k} \end{bmatrix}
\in \mathbb{R}^{2}
\label{eq:vector-estado}
\end{equation}

donde $V_{k}$ [\si{\metre\cubed}] representa el volumen de mineral retenido en
la cámara de trituración y $\dot{V}_{k}$ [\si{\metre\cubed\per\second}] denota
la tasa neta instantánea de variación de flujo volumétrico, dominada por la
extracción continua del chancador.

La dinámica expresa la conservación de masa de la tolva discretizada por Euler:

\begin{align}
V_{k+1} &= V_{k} + \Delta t\,\dot{V}_{k} + \Delta t\,u_{k}\\
\dot{V}_{k+1} &= \dot{V}_{k} + w_{k}\\
y_{k} &= V_{k} + v_{k}
\end{align}

lo que corresponde a las matrices

\begin{equation}
\mathbf{A} = \begin{bmatrix} 1 & \Delta t\\ 0 & 1\end{bmatrix},\qquad
\mathbf{B} = \begin{bmatrix} \Delta t\\ 0\end{bmatrix},\qquad
\mathbf{H} = \begin{bmatrix} 1 & 0\end{bmatrix}
\label{eq:matrices}
\end{equation}

La entrada $u_{k}$ [\si{\metre\cubed\per\second}] corresponde al caudal de
descarga de camión: es un evento operacional conocido y registrado por el SCADA,
y por eso ingresa como entrada de control y no como perturbación. En cambio
$\dot{V}_{k}$ representa la componente \textbf{no sensada} de la dinámica —la
tasa de extracción del chancador, que varía con la dureza del mineral y no
cuenta con instrumentación directa— y por eso se modela como caminata aleatoria.

Esta separación es el núcleo formal del problema de observabilidad parcial: la
estructura de $\mathbf{H}$ muestra que el sensor observa el volumen y nunca la
tasa.

\subsection{Verificación de la condición de observabilidad}

La reconstrucción de $\dot{V}_{k}$ a partir de la secuencia de mediciones de
volumen solo es posible si el par $(\mathbf{A}, \mathbf{H})$ es observable. La
matriz de observabilidad del sistema es

\begin{equation}
\mathcal{O} = \begin{bmatrix} \mathbf{H} \\ \mathbf{H}\mathbf{A}\end{bmatrix}
= \begin{bmatrix} 1 & 0\\ 1 & \Delta t\end{bmatrix},
\qquad \det(\mathcal{O}) = \Delta t \neq 0
\label{eq:observabilidad}
\end{equation}

de modo que $\operatorname{rango}(\mathcal{O}) = 2$ para todo paso de
discretización no nulo. El sistema es, por tanto, completamente observable: la
tasa de llenado es matemáticamente reconstruible pese a no medirse jamás de
forma directa. Esta verificación se ejecuta computacionalmente al inicio de cada
corrida del módulo y su resultado se registra en la salida del programa.

\subsection{Parametrización y criterios de selección}

La tabla~\ref{tab:parametros} consolida la parametrización completa del
experimento con sus unidades y su fundamento físico.

\begin{table}[H]
\centering
\caption{Parametrización del experimento y fundamento de cada valor.}
\label{tab:parametros}
\footnotesize
\begin{tabular}{@{}llll@{}}
\toprule
\textbf{Parámetro} & \textbf{Valor} & \textbf{Unidad} & \textbf{Fundamento}\\
\midrule
$\Delta t$ & \num{1,0} & \si{\second} & Muestreo SCADA a \SI{1}{\hertz}\\
$N$ & \num{1000} & pasos & Horizonte de \SI{16,7}{\minute} de operación\\
$M$ & \num{50} & corridas & Diseño de Monte Carlo\\
Semilla & \num{42} & --- & Reproducibilidad determinista\\
$Q_{11}$ & \num{0,01} & \si{\metre\tothe{6}} & Ruido de proceso sobre el volumen\\
$Q_{22}$ & \num{1,0e-4} & \si{\metre\tothe{6}\per\second\squared} &
Variabilidad de la tasa por dureza del mineral\\
$R$ & \num{5,0} & \si{\metre\tothe{6}} & Ruido del sensor
($\sigma_v = \SI{2,236}{\metre\cubed}$)\\
$V_{0}$ & \num{15,0} & \si{\metre\cubed} & Tolva casi vacía al inicio de turno\\
$\dot{V}_{0}$ & \num{-0,8} & \si{\metre\cubed\per\second} &
Extracción nominal ($\approx\SI{5760}{\tonne\per\hour}$ a
\SI{2,0}{\tonne\per\metre\cubed})\\
$\hat{V}_{0}$ & \num{15,0} & \si{\metre\cubed} & Estimación inicial del filtro\\
$\hat{\dot{V}}_{0}$ & \num{0,0} & \si{\metre\cubed\per\second} &
El filtro desconoce la tasa: debe inferirla\\
$\mathbf{P}_{0}$ & $\operatorname{diag}(10{,}0;\ 1{,}0)$ & --- &
Incertidumbre inicial\\
$u$ camión & \num{4,5} & \si{\metre\cubed\per\second} &
$\approx\SI{202}{\metre\cubed}$ por evento de \SI{45}{\second}\\
\bottomrule
\end{tabular}
\end{table}

La calibración física del régimen operacional responde a la escala real del
activo. Una extracción nominal de \SI{0,8}{\metre\cubed\per\second} equivale a
aproximadamente \SI{5760}{\tonne\per\hour} con densidad aparente de
\SI{2,0}{\tonne\per\metre\cubed}, coherente con un chancador giratorio de gran
capacidad. Cada descarga de camión aporta unos \SI{202}{\metre\cubed},
consistente con un camión de extracción de \SIrange{300}{400}{\tonne}. Cuatro
ciclos de camión en el horizonte simulado mantienen la tolva en una banda
operativa comprendida entre \num{7} y \SI{300}{\metre\cubed}.

\subsubsection*{Criterios de inclusión y exclusión}

Se incluyen regímenes de operación con pulsos de carga volumétrica en el rango
$u_{k} \in [0;\ 4{,}5]$~\si{\metre\cubed\per\second}, representativos del ciclo
de descarga de camiones acoplado a una tasa de trituración continua. Se excluyen
los estados transitorios de arranque y parada del equipo, así como las anomalías
catastróficas de atascamiento, por corresponder a regímenes en que el supuesto
de linealidad del modelo deja de ser válido.

\subsection{Nota de recalibración de la covarianza de proceso}
\label{sec:recalibracion}

Las entregas de la Fase 3 declaraban $\mathbf{Q} = \operatorname{diag}(0{,}01;\
0{,}05)$. Al implementar el modelo se verificó numéricamente que una covarianza
de proceso de $0{,}05$ por paso sobre la tasa produce, en un horizonte de $N =
1000$, una caminata aleatoria con desviación acumulada de
$\sqrt{1000 \times 0{,}05} = \SI{7,07}{\metre\cubed\per\second}$, es decir casi
nueve veces la tasa de extracción nominal de
\SI{0,8}{\metre\cubed\per\second}. Bajo ese régimen la tolva se llenaría o
vaciaría a velocidades físicamente imposibles para el activo, y el
\emph{truth model} dejaría de representar el proceso que pretende modelar.

Se recalibró en consecuencia $Q_{22} = 1{,}0\times10^{-4}$, valor que produce
una dispersión acumulada de \SI{0,32}{\metre\cubed\per\second} en torno a la
tasa nominal —aproximadamente un \SI{40}{\percent} de variabilidad atribuible a
la dureza variable del mineral—, coherente con la operación real de una estación
de chancado. El valor $Q_{11} = 0{,}01$ se mantuvo sin modificación.

Asimismo, se unificó la condición inicial en
$\mathbf{x}_{0} = [15{,}0;\ -0{,}8]^{\mathsf{T}}$, resolviendo una
inconsistencia entre documentos previos que declaraban alternativamente
$[15{,}0;\ 0{,}0]^{\mathsf{T}}$ y $[10{,}0;\ 0{,}0]^{\mathsf{T}}$. La
componente de tasa se fija en el valor nominal de extracción, mientras que la
estimación inicial del filtro se mantiene en $[15{,}0;\ 0{,}0]^{\mathsf{T}}$
precisamente para que el estimador deba inferir la tasa sin conocerla, poniendo
a prueba la propiedad de observabilidad verificada en la
ecuación~\eqref{eq:observabilidad}.
